\documentclass[12pt]{article}

% ============================================================
% Packages
% ============================================================
\usepackage{amsmath, amssymb, amsthm}
\usepackage{mathtools}
\usepackage{booktabs}
\usepackage{hyperref}
\usepackage[utf8]{inputenc}

% ============================================================
% Theorem environments
% ============================================================
\newtheorem{theorem}{Theorem}
\newtheorem{property}[theorem]{Structural Property}
\theoremstyle{remark}
\newtheorem{remark}[theorem]{Remark}

% ============================================================
% Notation (shared with main)
% ============================================================
\newcommand{\RR}{\mathbb{R}}
\newcommand{\ZZ}{\mathbb{Z}}
\newcommand{\CC}{\mathbb{C}}
\newcommand{\kk}{\mathbf{k}}
\newcommand{\LL}{\mathbf{L}}
\newcommand{\nn}{\mathbf{n}}
\newcommand{\Vol}{\mathrm{Vol}}
\DeclareMathOperator{\tr}{tr}
\DeclareMathOperator{\im}{im}
\DeclareMathOperator{\cker}{ker}

% ============================================================
% Title
% ============================================================
\title{Supplementary Material:\\
Exactness-preserving discrete de~Rham complexes\\
under Bloch-periodic boundary conditions}
\author{Alexandru Toader\\
Independent researcher, Buz\u{a}u, Romania\\
\texttt{toader\_alexandru@yahoo.com}}
\date{}

\begin{document}
\maketitle

% ============================================================
\section{Proof of Voronoi metric isotropy}
\label{sec:S1}
% ============================================================

\begin{property}[Voronoi isotropy]
On a periodic Voronoi tessellation of~$T^3$ generated by a point set in
general position (no four generators coplanar, no five cocircular), the
discrete metric tensors satisfy $G = \Vol \cdot I$ and $H = \Vol \cdot I$.
\end{property}

\paragraph{Proof ($G = \Vol \cdot I$)}
The edge metric tensor is
\[
  G_{ij} = \sum_e \star_1[e]\,(\Delta x_e)_i\,(\Delta x_e)_j.
\]
On a Voronoi tessellation, each dual face~$\sigma^*_e$ is perpendicular to
primal edge~$e$, and $\star_1[e] = |\sigma^*_e|/|\sigma_e|$. Rewriting the
sum over dual cells (Delaunay simplices), each cell contributes
\[
  \sum_{e \in \text{cell}} \frac{|\sigma^*_e|}{|\sigma_e|}
  \,(\Delta x_e)_i\,(\Delta x_e)_j.
\]
Apply the divergence theorem to the vector field $f(\mathbf{x}) = x_i\,\hat{e}_j$
on each dual cell:
$\int_{\partial\text{cell}} x_i\,(\hat{e}_j \cdot \hat{n})\,dA =
\Vol_{\text{cell}} \cdot \delta_{ij}$.
On a Voronoi dual, the boundary flux through each dual face equals
$\frac{|\sigma^*_e|}{|\sigma_e|}(\Delta x_e)_i(\Delta x_e)_j$ because the
face normal is parallel to~$\Delta x_e$ and the face area is~$|\sigma^*_e|$.
Summing over all dual cells: $G_{ij} = \Vol \cdot \delta_{ij}$. \qed

\paragraph{Proof ($H = \Vol \cdot I$)}
The face metric tensor is
\[
  H_{ij} = \sum_f \star_2[f]\,(A_f)_i\,(A_f)_j
\]
where $A_f$ is the face area vector. Apply the divergence theorem to
$f(\mathbf{x}) = x_i\,\hat{e}_j$ on each primal Voronoi cell. The boundary
flux through face~$f$ equals $(A_f)_i\,(\hat{e}_j \cdot \hat{n}_f)\,|A_f|$.
On a Voronoi tessellation, $\star_2[f] = \text{dual edge length} / |A_f|$,
and $A_f = |A_f|\,\hat{n}_f$, so
$\star_2[f]\,(A_f)_i\,(A_f)_j$ captures the boundary contribution.
Summing over cells: $H_{ij} = \Vol \cdot \delta_{ij}$. \qed

\paragraph{Robustness}
The identity holds regardless of cell regularity---only the Voronoi property
(edge~$\perp$ dual face) is used. Near-degenerate configurations with
$\min(\star_1) = 3.5 \times 10^{-5}$ still satisfy
$\|G/\Vol - I\| = 3.5 \times 10^{-12}$ (test~T6.8). The general position
assumption is needed for the mesh builder, not for the identity.

% ============================================================
\section{Schur complement derivation of $\omega^2 = |\kk|^2 + O(|\kk|^4)$}
\label{sec:S2}
% ============================================================

The acoustic eigenvalue is the minimum of~$K(\kk)$ on the physical subspace
(orthogonal to $\cker K = \im d_0$). At small~$|\kk|$, Schur complement
reduction onto the 3D harmonic subspace at~$\Gamma$ gives a $3 \times 3$
effective operator.

Exactness ensures this subspace is well-defined (Proposition~2 of main text).
Metric isotropy $G = H = \Vol \cdot I$ makes the effective operator
proportional to $|\kk|^2 I$ at leading order. The discrete curl identity
\[
  \frac{\partial}{\partial k_\beta}
  \bigl[h_2^\dagger\,\star_2\,d_1\,h_\alpha\bigr]_{\kk=\mathbf{0}}
  = i\,\varepsilon_{\beta\gamma\alpha}
\]
is the discrete analogue of the continuum identity
$\partial_\beta A_\alpha - \partial_\alpha A_\beta =
\varepsilon_{\beta\alpha\gamma}\,B_\gamma$ relating curl to the Levi-Civita
tensor. It arises because the harmonic 1-forms at~$\Gamma$ (the three
cohomology generators of~$H^1(T^3)$) deform into curl-like operators at
small~$|\kk|$. This identity is verified numerically to~$10^{-11}$ on 7~structures.
An analytical derivation is deferred to~[Paper~A2]; here it is used
as an empirically established fact.

The $O(|\kk|^4)$ correction arises from lattice dispersion.

Both metric tensors contribute to~$c^2$: $\star_2$ appears in
$K = d_1^\dagger \star_2\,d_1$ (numerator), $\star_1$ appears in~$M$
(denominator). Perturbing~$\star_2$ by 10\% gives $\Delta c^2 = 0.024$;
perturbing~$\star_1$ gives $\Delta c^2 = 0.012$. The face metric dominates
because it enters the operator directly rather than through the mass matrix.

Full proof including converse ($c^2 = 1$ implies metric isotropy) and moment
conservation $\tr(K^n)$ for $n < \text{girth}$ will appear in~[Paper~A2].

\end{document}
